%% ma: L10-B13-0065
%% lop: 10
%% chuong: 5
%% bai: 13
%% dang: 10.13.13
%% loai: TL
%% muc: TH
%% dap_an: "-"
%% nguon: "chuyen-de-cac-so-dac-trung-cua-mau-so-lieu-khong-ghep-nhom-toan-10-chuong-trinh-moi.pdf (GV. Phan Nhật Linh) (P03), Bài 02 Dạng 2 Bài tập tự luận 2"
%% kien_thuc: "Lớp 10 Bài 13 (số trung bình, trung vị, tứ phân vị, mốt)"
%% trang_thai: chinh-thuc
%% ghi_chu: "Lời giải gốc phần Q3 chép nhầm số liệu của bài khác, đã viết lại. Tự duyệt theo yêu cầu giáo viên 22/10/2026."
\begin{ex}
Tấm lấy ngẫu nhiên $5$ hạt từ một mâm trộn lẫn hạt gạo và hạt thóc, đếm xem có bao nhiêu hạt gạo trong $5$ hạt đó rồi trả lại vào mâm. Lặp lại việc làm trên $1000$ lần, Tấm ghi lại kết quả ở bảng sau:
\begin{center}
\begin{tabular}{|c|c|c|c|c|c|c|}\hline
Số hạt gạo & $0$ & $1$ & $2$ & $3$ & $4$ & $5$\\\hline
Số lần & $78$ & $259$ & $346$ & $230$ & $77$ & $10$\\\hline
\end{tabular}
\end{center}
Hãy tìm số trung bình, các tứ phân vị và mốt của mẫu số liệu.
\loigiai{$\overline x=\dfrac{0\cdot78+1\cdot259+2\cdot346+3\cdot230+4\cdot77+5\cdot10}{1000}=1{,}999$.

Tần số tích lũy: $78$; $337$; $683$; $913$; $990$; $1000$.

$Q_2=\dfrac{x_{500}+x_{501}}2=2$; $Q_1=\dfrac{x_{250}+x_{251}}2=1$; $Q_3=\dfrac{x_{750}+x_{751}}2=3$.

Mốt $M_o=2$.}
\end{ex}
